\scanpage{74}
% archive foliation: 73

Supposons maintenant la condition suiv. satisfaite
(en plus de : \uline{tout autom de $E$ est intérieur})

(C) Pour tout $O\subset E$, $O\in\mathcal O$, la loi de composition
dans $\operatorname{Hom}_{\underline B_1}(O,O)=$ ensemble des
\inserted{$(O,O)$-}biidéaux spéciaux
\deleted{\unsure{pour $O$}} \deleted{est}
$(L,M\mapsto L.M)$ est commutative :
\[
LM=ML.
\]
\margin{Il suffit (et il faut si $A$ est \unsure{dedek.|noeth.} de dim $1$)
que ce soit vrai pour les localisés \deleted{\unsure{tous les}}
idéaux maximaux (\unsure{ou leurs complétions, ou hensél.}).
Dans le cas local, \deleted{\unsure{est vérifiée}}
\unsure{est équivalent à la relation}
$LM=ML$ pour $L,M$ des biidéaux \unsure{bilatères biprincipaux} de $O$.}

$C\in\operatorname{Ob}\underline B_0$ une catégorie de Deuring,
et $X\in\operatorname{Ob}C$, de sorte que
$O\simeq\operatorname{End}_C(X)\in\mathcal O$ ; on peut
\deleted{\unsure{supposer}}\inserted{choisir}
$E\simeq O_K$. \inserted{\unsure{cet isomorphisme}}
On va définir sur la gr-catégorie
\[
\uline{\operatorname{Eq}}(C)=
\operatorname{Hom}_{\underline B_0}(C,C)
\]
une \deleted{loi de} contrainte de commutativité,
\margin{\unsure{en fera une cat. de} \uuline{Picard stricte}.
\deleted{\unsure{par transport}}}
à partir du choix de $X$. En effet, on a une gr-équivalence
\[
\uline{\operatorname{Hom}}_{\underline B_0}(C,C)
\overset{\alpha_X}{\simeq}
\uline{\operatorname{Hom}}_{\underline B_1}(\unsure{C,C|O,O})
=\text{catégorie des }
\deleted{\unsure{biidéaux}}\inserted{bimodules}
\deleted{\unsure{spéciaux sur $(O,O)$}}\text{ spéciaux}
\]
cette dernière catégorie étant \inserted{cat. de Picard}
stricte comm. (en fait, les contraintes d'associativité \uline{et} de commutativité
y sont des identités). On peut donc définir une contrainte de comm.
stricte sur $\operatorname{Eq}(C)$ par transport de structure.
On va l'expliciter ainsi : si $L,M$ sont des $(O,O)$-bimodules spéciaux,
on définit
\[
c:L\otimes_O M\xrightarrow{\sim}M\otimes_O L
\]
en \deleted{\unsure{choisissant}} utilisant un isom de $(O,O)$ bimodules
\[
\begin{aligned}
L&\simeq L'\subset E\\
M&\simeq M'\subset E
\end{aligned}
\]
avec $L',M'$ des $(O,O)$-biidéaux, et notant que
\diagram{
\begin{array}{rcl}
L\otimes M&\simeq&L'M'\\
&&\Vert\\[-2pt]
M\otimes L&\simeq&M'L'
\end{array}
}
